\chapter{§1.7 百分数}
\label{sec:1.7}


\prerequisite{§1.5 分数，§1.6 小数}

\gradelevel{6 年级}


\section*{百分数的意义}


\begin{explore}


期末考试中，语文的优秀率是 $85\%$ ，数学的优秀率是 $\frac{4}{5}$ ，英语的优秀率是 $0.82$ 。哪个科目的优秀率最高？三个数的表示方式不同，直接比较不太方便——有没有统一的方式？

在日常生活中，我们经常看到这样的数字：

\begin{itemize}
  \item "今日降水概率 $60\%$ "
  \item "牛奶含钙量比去年提高了 $15\%$ "
  \item "某商品打八折，即优惠 $20\%$ "
\end{itemize}


这些 $\%$ 是什么意思？

\end{explore}

\begin{understand}


\textbf{百分数}表示一个数是另一个数的百分之几，用符号 $\%$ 表示。

\[
85\% = \frac{85}{100}
\]


读作"百分之八十五"。

百分数是一种特殊的分数——分母固定为 $100$ 的分数。它的好处是让不同的分数和小数统一到同一个"尺度"上，便于比较。

回到引入情境：
\begin{itemize}
  \item 语文优秀率 $= 85\%$
  \item 数学优秀率 $= \frac{4}{5} = \frac{80}{100} = 80\%$
  \item 英语优秀率 $= 0.82 = \frac{82}{100} = 82\%$
\end{itemize}


统一成百分数后， $85\% > 82\% > 80\%$ ，语文优秀率最高。

\textbf{百分数与分数的区别}：百分数只表示两个量的比率关系，不能带单位。比如可以说"出勤率是 $95\%$ "，但不能说"一条绳子长 $95\%$ 米"。分数既可以表示比率，也可以表示具体的量。

\end{understand}

\begin{workedexamples}


\textbf{例 1.} 班级有 $50$ 名同学，今天到了 $48$ 名。今天的出勤率是多少？

\textbf{解：} 出勤率 $= \frac{\text{到校人数}}{\text{总人数}} = \frac{48}{50} = \frac{96}{100} = 96\%$ 。

\textbf{例 2.} 一个水池容量 $200$ 吨，现在存水 $170$ 吨。水池的蓄水率是多少？还需注入多少吨水才能注满？

\textbf{解：}
\begin{itemize}
  \item 蓄水率 $= \frac{170}{200} = \frac{85}{100} = 85\%$ 。
  \item 还需注入 $200 - 170 = 30$ （吨）。
\end{itemize}


\textbf{例 3.} 小明参加投篮练习，投了 $25$ 次，投中了 $18$ 次。他的命中率是多少？

\textbf{解：} 命中率 $= \frac{18}{25} = \frac{72}{100} = 72\%$ 。

\end{workedexamples}

\begin{keytakeaway}


\begin{itemize}
  \item 百分数表示"一个数是另一个数的百分之几"，符号是 $\%$ 。
  \item 百分数的本质是分母为 $100$ 的分数。
  \item 百分数只表示比率关系，不表示具体的量，不带单位。
  \item 百分数的好处：统一尺度，便于比较。
\end{itemize}


\end{keytakeaway}

\begin{practice}


\begin{enumerate}
  \item 班级 $40$ 人中有 $15$ 人是女生。女生占全班的百分之几？
  \item 工厂生产了 $500$ 个零件，其中 $8$ 个不合格。合格率是多少？
  \item 一块菜地面积 $800$ 平方米，其中 $200$ 平方米种了黄瓜。种黄瓜的面积占总面积的百分之几？
  \item 某次考试满分 $120$ 分，小亮得了 $102$ 分。他得到了满分的百分之几？
  \item 图书馆有 $2000$ 本书，借出了 $1300$ 本。借出率是多少？剩余率是多少？
\end{enumerate}


\end{practice}

\section*{百分数与分数、小数的互化}


\begin{explore}


在不同的场合，同一个数量关系可能用分数、小数或百分数来表示。例如"半价"可以说是 $\frac{1}{2}$ 、 $0.5$ 或 $50\%$ 。在解题过程中，我们经常需要在它们之间互相转化。怎么转化最方便？

\end{explore}

\begin{understand}


\textbf{百分数化小数}：去掉 $\%$ ，小数点向左移两位。

\[
75\% = 0.75, \quad 8\% = 0.08, \quad 120\% = 1.20 = 1.2
\]


\textbf{小数化百分数}：小数点向右移两位，添上 $\%$ 。

\[
0.45 = 45\%, \quad 1.2 = 120\%, \quad 0.036 = 3.6\%
\]


\textbf{百分数化分数}：先写成分母为 $100$ 的分数，再约分。

\[
75\% = \frac{75}{100} = \frac{3}{4}, \quad 12.5\% = \frac{12.5}{100} = \frac{125}{1000} = \frac{1}{8}
\]


\textbf{分数化百分数}：先用分子除以分母化成小数，再转为百分数。

\[
\frac{3}{8} = 3 \div 8 = 0.375 = 37.5\%
\]


\textbf{常用互化表}（建议记住）：

\begin{center}
\begin{tabular}{lll}
\toprule
分数 & 小数 & 百分数 \\
\midrule
$\frac{1}{2}$ & $0.5$ & $50\%$ \\
$\frac{1}{4}$ & $0.25$ & $25\%$ \\
$\frac{3}{4}$ & $0.75$ & $75\%$ \\
$\frac{1}{5}$ & $0.2$ & $20\%$ \\
$\frac{1}{8}$ & $0.125$ & $12.5\%$ \\
$\frac{1}{3}$ & $0.333\ldots$ & $33.\overline{3}\%$ \\
\bottomrule
\end{tabular}
\end{center}


\end{understand}

\begin{workedexamples}


\textbf{例 1.} 把 $\frac{5}{8}$ 化成百分数。

\textbf{解：} $\frac{5}{8} = 5 \div 8 = 0.625 = 62.5\%$ 。

\textbf{例 2.} 把 $160\%$ 化成分数和小数。

\textbf{解：}
\begin{itemize}
  \item 化小数： $160\% = 1.6$ 。
  \item 化分数： $160\% = \frac{160}{100} = \frac{8}{5} = 1\frac{3}{5}$ 。
\end{itemize}


\textbf{例 3.} 把下面三个数从小到大排列： $\frac{7}{20}$ ， $0.36$ ， $34\%$ 。

\textbf{解：} 统一化为小数比较：
\begin{itemize}
  \item $\frac{7}{20} = 0.35$
  \item $0.36 = 0.36$
  \item $34\% = 0.34$
\end{itemize}


从小到大： $34\% < \frac{7}{20} < 0.36$ 。

\textbf{例 4.} 甲数是 $42$ ，乙数是 $56$ 。甲数是乙数的百分之几？

\textbf{解：} $42 \div 56 = 0.75 = 75\%$ 。甲数是乙数的 $75\%$ 。

\end{workedexamples}

\begin{keytakeaway}


\begin{itemize}
  \item 百分数、小数、分数之间可以互化，核心桥梁是"除以 $100$ "与"乘以 $100$ "。
  \item 百分数化小数：去 $\%$ ，左移两位。小数化百分数：右移两位，加 $\%$ 。
  \item 分数化百分数：先做除法化为小数。
  \item 常见分数对应的百分数建议记熟。
\end{itemize}


\end{keytakeaway}

\begin{practice}


\begin{enumerate}
  \item 把 $45\%$ 化成小数和最简分数。
  \item 把 $\frac{5}{6}$ 化成百分数（保留一位小数）。
  \item 把 $0.08$ 化成百分数和分数。
  \item 按从小到大排列： $\frac{2}{5}$ ， $39\%$ ， $0.41$ 。
  \item 把 $125\%$ 化成小数和分数。
\end{enumerate}


\end{practice}

\section*{百分数的应用}


\begin{explore}


国庆节期间，商场搞促销活动。一件原价 $200$ 元的外套打八折出售。"打八折"是什么意思？折后价是多少？如果小红的妈妈在银行存了 $10000$ 元，年利率 $2.5\%$ ，一年后利息是多少？

百分数在日常生活中的应用非常广泛。这一节我们重点学习四类常见问题。

\end{explore}

\begin{understand}


\textbf{折扣问题}

"打八折"就是按原价的 $80\%$ 出售。"打七五折"就是按原价的 $75\%$ 出售。

\[
\text{折后价} = \text{原价} \times \text{折扣百分数}
\]


\[
\text{优惠金额} = \text{原价} - \text{折后价} = \text{原价} \times (1 - \text{折扣百分数})
\]


\textbf{利率问题}

把钱存进银行，银行会按一定的比率付给你利息。

\[
\text{利息} = \text{本金} \times \text{利率} \times \text{时间}
\]


例如，本金 $10000$ 元，年利率 $2.5\%$ ，存 $2$ 年：

\[
\text{利息} = 10000 \times 2.5\% \times 2 = 10000 \times 0.025 \times 2 = 500 \text{（元）}
\]


\textbf{增长率与减少率}

\[
\text{增长后的量} = \text{原量} \times (1 + \text{增长率})
\]


\[
\text{减少后的量} = \text{原量} \times (1 - \text{减少率})
\]


\textbf{浓度问题}

\[
\text{浓度} = \frac{\text{溶质质量}}{\text{溶液质量}} \times 100\%
\]


其中，溶液质量 $=$ 溶质质量 $+$ 溶剂质量。

\end{understand}

\begin{workedexamples}


\textbf{例 1.} 一双运动鞋原价 $350$ 元，现在打七折出售。折后价是多少元？便宜了多少元？

\textbf{解：}
\begin{itemize}
  \item 折后价 $= 350 \times 70\% = 350 \times 0.7 = 245$ （元）。
  \item 便宜了 $350 - 245 = 105$ （元）。
\end{itemize}


\textbf{例 2.} 爸爸在银行存了 $50000$ 元定期存款，年利率 $3.2\%$ ，存期 $3$ 年。到期后利息是多少？本金和利息一共是多少？

\textbf{解：}
\begin{itemize}
  \item 利息 $= 50000 \times 3.2\% \times 3 = 50000 \times 0.032 \times 3 = 4800$ （元）。
  \item 本金加利息 $= 50000 + 4800 = 54800$ （元）。
\end{itemize}


\textbf{例 3.} 某市去年人口 $800$ 万，今年比去年增长了 $1.5\%$ 。今年人口是多少万？

\textbf{解：}

\[
800 \times (1 + 1.5\%) = 800 \times 1.015 = 812 \text{（万人）}
\]


\textbf{例 4.} 用 $40$ 克盐溶解在 $160$ 克水中，盐水的浓度是多少？

\textbf{解：}
\begin{itemize}
  \item 溶液质量 $= 40 + 160 = 200$ （克）。
  \item 浓度 $= \frac{40}{200} \times 100\% = 20\%$ 。
\end{itemize}


\textbf{例 5.} 一件商品原价 $a$ 元，先涨价 $10\%$ ，再打九折（ $90\%$ ）出售。最终售价是多少？是否等于原价？

\textbf{解：}
\begin{itemize}
  \item 涨价后： $a \times (1 + 10\%) = 1.1a$ （元）。
  \item 打九折后： $1.1a \times 90\% = 1.1a \times 0.9 = 0.99a$ （元）。
\end{itemize}


最终售价是 $0.99a$ 元，比原价少了 $1\%$ ，并不等于原价。

\end{workedexamples}

\begin{quote}
这说明"先涨 $10\%$ 再降 $10\%$ "并不能还原为原价。这是百分数应用中一个常见的"陷阱"。
\end{quote}


\textbf{例 6.} 某商品打八五折后售价 $170$ 元，原价是多少？

\textbf{解：} 设原价为 $x$ 元。

\[
x \times 85\% = 170
\]


\[
x = 170 \div 0.85 = 200 \text{（元）}
\]


原价是 $200$ 元。

\begin{keytakeaway}


\begin{itemize}
  \item 折扣：打 $n$ 折就是按原价的 $n \times 10\%$ 出售。
  \item 利息 $=$ 本金 $\times$ 利率 $\times$ 时间。
  \item 增减问题的关键：在原量的基础上" $\times (1 \pm \text{百分率})$ "。
  \item "先增后减"或"先减后增"同一百分率，结果不等于原值。
  \item 浓度 $=$ 溶质 $\div$ 溶液 $\times 100\%$ 。
\end{itemize}


\end{keytakeaway}

\begin{practice}


\begin{enumerate}
  \item 一台电脑原价 $5000$ 元，打八折出售。折后价是多少？
  \item 妈妈存了 $20000$ 元，年利率 $2.8\%$ ，存 $2$ 年。到期后利息是多少？
  \item 某学校去年有学生 $1200$ 人，今年比去年减少了 $5\%$ 。今年有多少人？
  \item 用 $25$ 克糖溶解在 $475$ 克水中，糖水的浓度是多少？
  \item 一辆自行车打七五折后售价 $450$ 元，原价是多少？
  \item 小明家的水费上个月是 $80$ 元，这个月比上个月增加了 $15\%$ 。这个月水费是多少元？
  \item 某品牌牛奶标注"含钙量增加 $25\%$ "。如果原来每 $100$ 毫升含钙 $120$ 毫克，现在每 $100$ 毫升含钙多少毫克？
  \item 配制 $500$ 克浓度为 $16\%$ 的盐水，需要多少克盐和多少克水？
\end{enumerate}


\end{practice}



\begin{answer}


\textbf{练一练 百分数的意义}

\begin{enumerate}
  \item $\frac{15}{40} = \frac{3}{8} = 37.5\%$ 。
  \item 不合格率 $= \frac{8}{500} = 1.6\%$ ，合格率 $= 100\% - 1.6\% = 98.4\%$ 。
  \item $\frac{200}{800} = 25\%$ 。
  \item $\frac{102}{120} = 0.85 = 85\%$ 。
  \item 借出率 $= \frac{1300}{2000} = 65\%$ ，剩余率 $= 100\% - 65\% = 35\%$ 。
\end{enumerate}


\textbf{练一练 百分数与分数、小数的互化}

\begin{enumerate}
  \item $45\% = 0.45 = \frac{45}{100} = \frac{9}{20}$ 。
  \item $\frac{5}{6} = 5 \div 6 \approx 0.833 = 83.3\%$ 。
  \item $0.08 = 8\% = \frac{8}{100} = \frac{2}{25}$ 。
  \item $39\% = 0.39$ ， $\frac{2}{5} = 0.4$ ， $0.41 = 0.41$ 。从小到大： $39\% < \frac{2}{5} < 0.41$ 。
  \item $125\% = 1.25 = \frac{125}{100} = \frac{5}{4} = 1\frac{1}{4}$ 。
\end{enumerate}


\textbf{练一练 百分数的应用}

\begin{enumerate}
  \item $5000 \times 80\% = 5000 \times 0.8 = 4000$ （元）。
  \item $20000 \times 2.8\% \times 2 = 20000 \times 0.028 \times 2 = 1120$ （元）。
  \item $1200 \times (1 - 5\%) = 1200 \times 0.95 = 1140$ （人）。
  \item 溶液质量 $= 25 + 475 = 500$ （克），浓度 $= \frac{25}{500} \times 100\% = 5\%$ 。
  \item $450 \div 75\% = 450 \div 0.75 = 600$ （元）。
  \item $80 \times (1 + 15\%) = 80 \times 1.15 = 92$ （元）。
  \item $120 \times (1 + 25\%) = 120 \times 1.25 = 150$ （毫克）。
  \item 盐的质量 $= 500 \times 16\% = 80$ （克），水的质量 $= 500 - 80 = 420$ （克）。
\end{enumerate}


\end{answer}
